\documentclass[../main.tex]{subfiles}
\begin{document}

\chapter{Thread Performance Considerations}
\lessoninfo{
  video={P2L5},
  textbook={OSTEP \cite{ostep} ch.~26, 33; Silberschatz et al.\ \cite{silberschatz2018} ch.~4--5; Kerrisk \cite{kerrisk2010} ch.~49--50, 63; Pai et al.\ \cite{pai1999}},
  keywords={performance metric, execution time, throughput, response time, multi-process model, multithreaded model, event-driven model, event dispatcher, asynchronous I/O, helper, AMPED, Flash, workload, baseline},
}

Lessons~3--5 introduced threads, their programming interface and their
implementation, and argued that multithreading makes applications faster
and more efficient.  This lesson asks how such a claim is evaluated.  It
shows that which design is better depends on the performance metric and on
the workload, compares the multi-process, multithreaded and event-driven
architectures of a web server, and follows the evaluation of the Flash web
server of Pai et al.\ \cite{pai1999} as a case study in the design of
experiments.

\section{Comparing threading models}
\label{sec:l06-compare}

Lesson~3 compared the boss-workers and pipeline patterns by the total time
they need to complete a batch of requests.  That is only one of several
reasonable criteria, and a different criterion can reverse the
verdict.\source{P2L5, which threading model is better.}  Consider a server
with five threads that receives a batch of $n$ requests, each of which
needs \SI{100}{\milli\second} of processing:
\begin{itemize}
  \item in the boss-workers configuration one thread is the boss, whose own
    time per request is negligible, and $w = 4$ threads are workers, each of
    which processes a whole request in $t = \SI{100}{\milli\second}$;
  \item in the pipeline configuration each of the $k = 5$ threads is a
    stage, and every stage takes $s = \SI{20}{\milli\second}$.
\end{itemize}
With boss-workers the requests complete in groups of $w$, so request $i$
completes at time $\lceil i/w \rceil\, t$; in the pipeline request $i$
completes at $(k + i - 1)\, s$.  Two metrics can be derived from these
completion times.  The time to complete the whole batch is the completion
time of the last request.  The \emph{average completion time} of a request
is
\begin{equation}
  \overline{T}_{\mathrm{bw}}(n) = \frac{t}{n} \sum_{i=1}^{n}
    \Bigl\lceil \frac{i}{w} \Bigr\rceil ,
  \qquad
  \overline{T}_{\mathrm{pl}}(n) = k\, s + \frac{(n-1)\, s}{2} .
  \label{eq:l06-mean-completion}
\end{equation}

\begin{figure}[htbp]
  \centering
  % Completion times of nine requests: boss-workers with four workers
% (100 ms per request) and a five-stage pipeline (20 ms per stage).
% Usage: % Completion times of nine requests: boss-workers with four workers
% (100 ms per request) and a five-stage pipeline (20 ms per stage).
% Usage: % Completion times of nine requests: boss-workers with four workers
% (100 ms per request) and a five-stage pipeline (20 ms per stage).
% Usage: \input{figures/l06-completion}   (width ~116 mm)
\begin{tikzpicture}[x=1mm, y=1mm, font=\footnotesize\sffamily,
  >={Stealth[length=1.6mm]},
  lab/.style={font=\scriptsize\sffamily},
  row/.style={font=\scriptsize\sffamily, anchor=east},
  done/.style={draw, circle, fill=ln-accent!20, minimum size=4.2mm,
    inner sep=0pt, font=\scriptsize\sffamily},
  mean/.style={ln-caution, densely dashed, thick}]
  % x = 24 + 0.3 t, with t in ms
  % boss-workers
  \node[row] at (22,18) {\iflangja{ボス・ワーカー}{boss-workers}};
  \draw[ln-rule] (24,18) -- (114,18);
  \draw[mean] (74,15) -- (74,23);
  \node[lab, anchor=south, text=ln-caution] at (74,23.1)
    {\iflangja{平均 167\,ms}{mean 167\,ms}};
  \node[done] at (54,18) {4};
  \node[done] at (84,18) {4};
  \node[done] at (114,18) {1};
  % pipeline
  \node[row] at (22,4) {\iflangja{パイプライン}{pipeline}};
  \draw[ln-rule] (24,4) -- (114,4);
  \draw[mean] (78,-0.5) -- (78,8.5);
  \node[lab, anchor=south, text=ln-caution] at (78,8.6)
    {\iflangja{平均 180\,ms}{mean 180\,ms}};
  \foreach \t in {100,120,...,260} { \node[done] at (24+0.3*\t,4) {1}; }
  % time axis
  \draw[->] (24,-4) -- (116,-4);
  \foreach \t in {0,50,...,300} {
    \draw (24+0.3*\t,-4) -- (24+0.3*\t,-5);
    \node[lab, anchor=north] at (24+0.3*\t,-5.2) {\t};
  }
  \node[row] at (22,-4) {\iflangja{時間 (ms)}{time (ms)}};
\end{tikzpicture}
   (width ~116 mm)
\begin{tikzpicture}[x=1mm, y=1mm, font=\footnotesize\sffamily,
  >={Stealth[length=1.6mm]},
  lab/.style={font=\scriptsize\sffamily},
  row/.style={font=\scriptsize\sffamily, anchor=east},
  done/.style={draw, circle, fill=ln-accent!20, minimum size=4.2mm,
    inner sep=0pt, font=\scriptsize\sffamily},
  mean/.style={ln-caution, densely dashed, thick}]
  % x = 24 + 0.3 t, with t in ms
  % boss-workers
  \node[row] at (22,18) {\iflangja{ボス・ワーカー}{boss-workers}};
  \draw[ln-rule] (24,18) -- (114,18);
  \draw[mean] (74,15) -- (74,23);
  \node[lab, anchor=south, text=ln-caution] at (74,23.1)
    {\iflangja{平均 167\,ms}{mean 167\,ms}};
  \node[done] at (54,18) {4};
  \node[done] at (84,18) {4};
  \node[done] at (114,18) {1};
  % pipeline
  \node[row] at (22,4) {\iflangja{パイプライン}{pipeline}};
  \draw[ln-rule] (24,4) -- (114,4);
  \draw[mean] (78,-0.5) -- (78,8.5);
  \node[lab, anchor=south, text=ln-caution] at (78,8.6)
    {\iflangja{平均 180\,ms}{mean 180\,ms}};
  \foreach \t in {100,120,...,260} { \node[done] at (24+0.3*\t,4) {1}; }
  % time axis
  \draw[->] (24,-4) -- (116,-4);
  \foreach \t in {0,50,...,300} {
    \draw (24+0.3*\t,-4) -- (24+0.3*\t,-5);
    \node[lab, anchor=north] at (24+0.3*\t,-5.2) {\t};
  }
  \node[row] at (22,-4) {\iflangja{時間 (ms)}{time (ms)}};
\end{tikzpicture}
   (width ~116 mm)
\begin{tikzpicture}[x=1mm, y=1mm, font=\footnotesize\sffamily,
  >={Stealth[length=1.6mm]},
  lab/.style={font=\scriptsize\sffamily},
  row/.style={font=\scriptsize\sffamily, anchor=east},
  done/.style={draw, circle, fill=ln-accent!20, minimum size=4.2mm,
    inner sep=0pt, font=\scriptsize\sffamily},
  mean/.style={ln-caution, densely dashed, thick}]
  % x = 24 + 0.3 t, with t in ms
  % boss-workers
  \node[row] at (22,18) {\iflangja{ボス・ワーカー}{boss-workers}};
  \draw[ln-rule] (24,18) -- (114,18);
  \draw[mean] (74,15) -- (74,23);
  \node[lab, anchor=south, text=ln-caution] at (74,23.1)
    {\iflangja{平均 167\,ms}{mean 167\,ms}};
  \node[done] at (54,18) {4};
  \node[done] at (84,18) {4};
  \node[done] at (114,18) {1};
  % pipeline
  \node[row] at (22,4) {\iflangja{パイプライン}{pipeline}};
  \draw[ln-rule] (24,4) -- (114,4);
  \draw[mean] (78,-0.5) -- (78,8.5);
  \node[lab, anchor=south, text=ln-caution] at (78,8.6)
    {\iflangja{平均 180\,ms}{mean 180\,ms}};
  \foreach \t in {100,120,...,260} { \node[done] at (24+0.3*\t,4) {1}; }
  % time axis
  \draw[->] (24,-4) -- (116,-4);
  \foreach \t in {0,50,...,300} {
    \draw (24+0.3*\t,-4) -- (24+0.3*\t,-5);
    \node[lab, anchor=north] at (24+0.3*\t,-5.2) {\t};
  }
  \node[row] at (22,-4) {\iflangja{時間 (ms)}{time (ms)}};
\end{tikzpicture}

  \caption{Completion times of nine requests.  Each circle gives the number
    of requests completing at that time: boss-workers with four workers
    completes four requests every \SI{100}{\milli\second}, the five-stage
    pipeline one request every \SI{20}{\milli\second} after the first
    \SI{100}{\milli\second}.  Dashed lines mark the average completion
    times.}
  \label{fig:l06-completion}
\end{figure}

\begin{example}
  \label{ex:l06-completion}
  For a batch of $n = 9$ requests (\cref{fig:l06-completion}), boss-workers
  completes four requests at \SI{100}{\milli\second}, four at
  \SI{200}{\milli\second} and the last at \SI{300}{\milli\second}; the
  pipeline completes the first request at \SI{100}{\milli\second} and then
  one every \SI{20}{\milli\second}, the last at \SI{260}{\milli\second}.
  The pipeline finishes the batch sooner.  By
  \cref{eq:l06-mean-completion}, however, the average completion time is
  $(4 \cdot 100 + 4 \cdot 200 + 300)/9 \approx \SI{166.7}{\milli\second}$
  for boss-workers and $100 + 8 \cdot 20/2 = \SI{180}{\milli\second}$ for
  the pipeline, so boss-workers serves the average request sooner.  For
  $n = 40$ the pipeline wins by both metrics: it finishes the batch at
  \SI{880}{\milli\second} rather than \SI{1000}{\milli\second}, with an
  average of \SI{490}{\milli\second} rather than \SI{550}{\milli\second}.
\end{example}

Boss-workers completes $w$ requests at once after $t$, so many requests
finish early, whereas the pipeline delivers only one request per stage
time.  In steady state, however, the pipeline completes $1/s = 50$ requests
per second against $w/t = 40$ for boss-workers, and in a large batch this
higher rate dominates.  Which model is better thus depends both on the
metric and on the number of requests.

\needspace{6\baselineskip}
\section{Usefulness of threads}
\label{sec:l06-useful}

The earlier lessons gave several reasons for using
threads.\source{P2L5, are threads useful, visual metaphor; OSTEP ch.~26.}
\begin{itemize}
  \item \textbf{Parallelization}: on several CPUs, threads that work on
    parts of the same task speed it up.
  \item \textbf{Specialization}: threads dedicated to particular tasks work
    with a hot cache and can be given different priorities.
  \item \textbf{Efficiency}: threads share an address space, so they need
    less memory than processes and communicate and synchronize more cheaply.
  \item \textbf{Hiding idle time}: even on a single CPU, another thread can
    run while one thread waits for I/O.
\end{itemize}
Whether these benefits make a thread-based design useful depends on what is
measured, and different parties measure different things.  For an
application that multiplies matrices, the relevant measure is its execution
time.  For a web service it is the number of client requests handled per
unit time, and the time a client waits for a response---which may be
reported as the average, the maximum or a percentile, such as the 95th
percentile.  For the owner of the hardware it is the utilization of
resources such as the CPU.  Three measures thus recur for any system that
serves requests: how many requests it completes per unit time, how long a
request takes to be answered, and how much of the capacity of the system is
in use.

\begin{example}
  A server handles 200 requests: 180 receive a response after
  \SI{20}{\milli\second}, and 20 after \SI{400}{\milli\second}.  The average
  response time is $(180 \cdot 20 + 20 \cdot 400)/200 =
  \SI{58}{\milli\second}$.  The 95th percentile is the smallest value that
  at least 95\,\% of the requests do not exceed; in sorted order the 190th
  response time is \SI{400}{\milli\second}, so the 95th percentile is
  \SI{400}{\milli\second}.  The average hides the fact that one request in
  ten waits twenty times longer than most.
\end{example}

\needspace{6\baselineskip}
\section{Performance metrics}
\label{sec:l06-metrics}

A statement that one design performs better than another is meaningful
only with respect to a metric.\source{P2L5, performance metrics; really
are threads useful; Silberschatz et al.\ ch.~5.}

\begin{definition}[Performance metric]
  A measurement standard for a system, that is, a measurable, quantifiable
  property of the system of interest that can be used to evaluate its
  behaviour, for instance to compare it with alternative implementations,
  is a \term{performance metric}[性能指標].
\end{definition}

Ideally a metric is obtained by experiments with real software deployed on
real machines under real workloads.  This is often infeasible---the
machines may not be available, the software may not be finished, and real
users cannot be asked to wait for an experiment---so measurements come
instead from small experiments designed to be representative of realistic
settings, or from simulation.  An environment of this kind, set up to run
experiments and obtain metrics, is called a \term{testbed}[テストベッド].

The most common metrics are the following.
\begin{description}
  \item[Execution time] The \term{execution time}[実行時間] of a task is the
    time from its start to its completion.  It is the natural metric for a
    batch computation such as a matrix multiplication.
  \item[Throughput] The \term{throughput}[スループット] of a system is the
    number of tasks, such as requests, that it completes per unit time.  It
    differs from the \emph{request rate}, the number of requests that
    arrive per unit time: a system that cannot keep up with the request rate
    still completes requests only at its throughput, and the backlog grows.
  \item[Response time] The \term{response time}[応答時間] of a request is
    the time from its issue until the client receives the response,
    including any time the request waits before it is processed.  Over many
    requests it is summarized by an average, a maximum or a percentile.
  \item[Wait time] The \term{wait time}[待ち時間] of a task is the time from
    its submission until it starts to execute.  It matters when users care
    about when their jobs start rather than when they finish.
  \item[CPU utilization] The \term{CPU utilization}[CPU使用率] is the
    fraction of time during which the CPU performs useful work; utilization
    metrics for other resources are defined in the same way.
\end{description}

Many further metrics relate performance to what it costs or to what a
particular party cares about (\cref{tab:l06-metrics}).  Providers of a
service often commit to performance targets in a
\term{service level agreement}[SLA]\index{SLA|see{service level agreement}}
(SLA), for example that 99\,\% of requests receive a response within
\SI{200}{\milli\second}; the percentage of requests that violate the
agreement is then itself a metric.

\begin{table}[htbp]
  \centering
  \caption{Performance metrics and the parties that are typically
    interested in them.}
  \label{tab:l06-metrics}
  \small
  \begin{tabular}{@{}>{\raggedright\arraybackslash}p{0.25\linewidth}>{\raggedright\arraybackslash}p{0.45\linewidth}>{\raggedright\arraybackslash}p{0.22\linewidth}@{}}
    \toprule
    Metric & Measures & Of interest to \\
    \midrule
    execution time & time to complete a task & users of batch jobs \\
    throughput, request rate & tasks completed, requests arriving per unit
      time & operators \\
    response time & time until a request is answered & clients \\
    wait time & time until a task starts & users \\
    CPU utilization & fraction of time the CPU is busy & operators \\
    platform efficiency & throughput relative to the resources used &
      data center operators \\
    performance per dollar, per watt & performance relative to cost or power
      & operators, buyers \\
    SLA violations & percentage of requests that miss the agreed target &
      providers, clients \\
    client-perceived performance & quality as seen by a client, e.g.\ the
      frame rate of a video & clients \\
    aggregate performance & a metric averaged over many tasks or clients &
      operators \\
    average resource usage & e.g.\ the memory used over time & operators \\
    \bottomrule
  \end{tabular}
\end{table}

\subsection{It depends}

The question whether threads are useful therefore has no single answer.
The answer depends on the metric, as \cref{sec:l06-compare} showed for
average and total completion time, and it depends on the workload: the
better of two threading models changed with the number of requests.  The
same holds for most design decisions in operating systems.  Different kinds
of graphs favour different shortest-path algorithms, different patterns of
file access favour different file system designs, and different request
loads favour different server architectures.

\begin{keypoint}
  No design is best in general.  A comparison is meaningful only for a
  stated metric and a stated workload, and the answer to \enquote{which is
  better?} is typically \enquote{it depends}.
\end{keypoint}

\needspace{6\baselineskip}
\section{Multi-process and multithreaded servers}
\label{sec:l06-mp-mt}

The remainder of the lesson compares ways of structuring a concurrent
server, using as the example a web server that serves static
files.\source{P2L5, multi-process vs.\ multithreaded; Silberschatz et
al.\ ch.~4.}  A client sends a request,
the server processes it in the steps of \cref{fig:l06-web-steps}, and it
responds by sending the file.  Some steps, such as parsing the request and
computing the response header, need only CPU time.  Others may block:
accepting a connection waits for a client, reading the request and sending
the data wait for the network, and finding and reading the file may wait
for the disk.

\begin{figure}[htbp]
  \centering
  % Processing steps of a simple web server for static content; the bars
% mark the steps that may wait for the network or for the disk.
% Usage: % Processing steps of a simple web server for static content; the bars
% mark the steps that may wait for the network or for the disk.
% Usage: % Processing steps of a simple web server for static content; the bars
% mark the steps that may wait for the network or for the disk.
% Usage: \input{figures/l06-web-steps}   (width ~116 mm)
\begin{tikzpicture}[x=1mm, y=1mm, font=\scriptsize\sffamily,
  st/.style={draw, rounded corners=1.5pt, fill=ln-box, align=center,
    minimum width=14.6mm, minimum height=9mm, inner sep=0.5pt},
  num/.style={font=\tiny\sffamily, text=ln-muted, anchor=south},
  net/.style={draw=none, fill=ln-accent!60},
  dsk/.style={draw=none, fill=ln-caution!60},
  >={Stealth[length=1.2mm]}]
  \node[st] (s1) at (7.3,0)   {\iflangja{接続の\\受け付け}{accept\\connection}};
  \node[st] (s2) at (24.2,0)  {\iflangja{要求の\\読み込み}{read\\request}};
  \node[st] (s3) at (41.1,0)  {\iflangja{要求の\\解析}{parse\\request}};
  \node[st] (s4) at (58.0,0)  {\iflangja{ファイル\\の検索}{find\\file}};
  \node[st] (s5) at (74.9,0)  {\iflangja{ヘッダ\\の計算}{compute\\header}};
  \node[st] (s6) at (91.8,0)  {\iflangja{ヘッダ\\の送信}{send\\header}};
  \node[st] (s7) at (108.7,0) {\iflangja{ファイル\\読込・送信}{read file,\\send data}};
  \foreach \i [evaluate=\i as \j using int(\i+1)] in {1,...,6} {
    \draw[->] (s\i.east) -- (s\j.west);
  }
  \foreach \i in {1,...,7} { \node[num] at (s\i.north) {\i}; }
  % may wait for the network
  \foreach \x in {7.3,24.2,91.8,108.7} {
    \fill[net] (\x-7.3,-6.2) rectangle (\x+7.3,-5.0);
  }
  % may wait for the disk
  \foreach \x in {58.0,108.7} {
    \fill[dsk] (\x-7.3,-8.2) rectangle (\x+7.3,-7.0);
  }
  % legend
  \fill[net] (22,-12.6) rectangle (27,-11.4);
  \node[anchor=west] at (28,-12) {\iflangja{ネットワークを待つ可能性}{may wait for the network}};
  \fill[dsk] (68,-12.6) rectangle (73,-11.4);
  \node[anchor=west] at (74,-12) {\iflangja{ディスクを待つ可能性}{may wait for the disk}};
\end{tikzpicture}
   (width ~116 mm)
\begin{tikzpicture}[x=1mm, y=1mm, font=\scriptsize\sffamily,
  st/.style={draw, rounded corners=1.5pt, fill=ln-box, align=center,
    minimum width=14.6mm, minimum height=9mm, inner sep=0.5pt},
  num/.style={font=\tiny\sffamily, text=ln-muted, anchor=south},
  net/.style={draw=none, fill=ln-accent!60},
  dsk/.style={draw=none, fill=ln-caution!60},
  >={Stealth[length=1.2mm]}]
  \node[st] (s1) at (7.3,0)   {\iflangja{接続の\\受け付け}{accept\\connection}};
  \node[st] (s2) at (24.2,0)  {\iflangja{要求の\\読み込み}{read\\request}};
  \node[st] (s3) at (41.1,0)  {\iflangja{要求の\\解析}{parse\\request}};
  \node[st] (s4) at (58.0,0)  {\iflangja{ファイル\\の検索}{find\\file}};
  \node[st] (s5) at (74.9,0)  {\iflangja{ヘッダ\\の計算}{compute\\header}};
  \node[st] (s6) at (91.8,0)  {\iflangja{ヘッダ\\の送信}{send\\header}};
  \node[st] (s7) at (108.7,0) {\iflangja{ファイル\\読込・送信}{read file,\\send data}};
  \foreach \i [evaluate=\i as \j using int(\i+1)] in {1,...,6} {
    \draw[->] (s\i.east) -- (s\j.west);
  }
  \foreach \i in {1,...,7} { \node[num] at (s\i.north) {\i}; }
  % may wait for the network
  \foreach \x in {7.3,24.2,91.8,108.7} {
    \fill[net] (\x-7.3,-6.2) rectangle (\x+7.3,-5.0);
  }
  % may wait for the disk
  \foreach \x in {58.0,108.7} {
    \fill[dsk] (\x-7.3,-8.2) rectangle (\x+7.3,-7.0);
  }
  % legend
  \fill[net] (22,-12.6) rectangle (27,-11.4);
  \node[anchor=west] at (28,-12) {\iflangja{ネットワークを待つ可能性}{may wait for the network}};
  \fill[dsk] (68,-12.6) rectangle (73,-11.4);
  \node[anchor=west] at (74,-12) {\iflangja{ディスクを待つ可能性}{may wait for the disk}};
\end{tikzpicture}
   (width ~116 mm)
\begin{tikzpicture}[x=1mm, y=1mm, font=\scriptsize\sffamily,
  st/.style={draw, rounded corners=1.5pt, fill=ln-box, align=center,
    minimum width=14.6mm, minimum height=9mm, inner sep=0.5pt},
  num/.style={font=\tiny\sffamily, text=ln-muted, anchor=south},
  net/.style={draw=none, fill=ln-accent!60},
  dsk/.style={draw=none, fill=ln-caution!60},
  >={Stealth[length=1.2mm]}]
  \node[st] (s1) at (7.3,0)   {\iflangja{接続の\\受け付け}{accept\\connection}};
  \node[st] (s2) at (24.2,0)  {\iflangja{要求の\\読み込み}{read\\request}};
  \node[st] (s3) at (41.1,0)  {\iflangja{要求の\\解析}{parse\\request}};
  \node[st] (s4) at (58.0,0)  {\iflangja{ファイル\\の検索}{find\\file}};
  \node[st] (s5) at (74.9,0)  {\iflangja{ヘッダ\\の計算}{compute\\header}};
  \node[st] (s6) at (91.8,0)  {\iflangja{ヘッダ\\の送信}{send\\header}};
  \node[st] (s7) at (108.7,0) {\iflangja{ファイル\\読込・送信}{read file,\\send data}};
  \foreach \i [evaluate=\i as \j using int(\i+1)] in {1,...,6} {
    \draw[->] (s\i.east) -- (s\j.west);
  }
  \foreach \i in {1,...,7} { \node[num] at (s\i.north) {\i}; }
  % may wait for the network
  \foreach \x in {7.3,24.2,91.8,108.7} {
    \fill[net] (\x-7.3,-6.2) rectangle (\x+7.3,-5.0);
  }
  % may wait for the disk
  \foreach \x in {58.0,108.7} {
    \fill[dsk] (\x-7.3,-8.2) rectangle (\x+7.3,-7.0);
  }
  % legend
  \fill[net] (22,-12.6) rectangle (27,-11.4);
  \node[anchor=west] at (28,-12) {\iflangja{ネットワークを待つ可能性}{may wait for the network}};
  \fill[dsk] (68,-12.6) rectangle (73,-11.4);
  \node[anchor=west] at (74,-12) {\iflangja{ディスクを待つ可能性}{may wait for the disk}};
\end{tikzpicture}

  \caption{Processing steps of a simple web server for static content.
    Bars below a step show which resource it may have to wait for.}
  \label{fig:l06-web-steps}
\end{figure}

A server that executes these steps for one request at a time serves a
single client and keeps all others waiting, even while it is itself
blocked.  The simplest remedy is to run several copies of it.

\begin{definition}[Multi-process model]
  In the \term{multi-process model}[マルチプロセスモデル] (MP) a server
  consists of several processes, each of which runs a single-threaded
  server that executes all the steps for one request at a time
  (\cref{fig:l06-mp-mt}a).
\end{definition}

The multi-process model is easy to program, since each process is an
ordinary sequential program.  Its costs follow from the separate address
spaces.  Many processes need much memory, which leaves less memory for
caching files.  Switching between processes is costly.  State that the
processes should share, such as a cache of file data or of response
headers, is hard and expensive to maintain, because it requires IPC or
explicitly shared memory.  Finally, all processes must accept connections on
the same port, which requires some care in setting them up; typically a
parent process creates the listening socket and then forks the server
processes, which inherit the socket and all call \texttt{accept} on it.

\begin{definition}[Multithreaded model]
  In the \term{multithreaded model}[マルチスレッドモデル] (MT) a server is a
  single process whose threads, sharing one address space, each process
  requests concurrently (\cref{fig:l06-mp-mt}b).
\end{definition}

The threads can be organized in several ways: every thread may execute all
the steps, accepting connections itself; or, in the boss-workers pattern of
Lesson~3, a boss thread accepts connections and passes them to workers; or
the steps may be the stages of a pipeline.  In every organization the
threads share the address space, so shared state such as a cache is
directly accessible, the memory footprint is smaller, and context switches
between threads are cheap.  In return, the implementation is not simple: the
shared state requires synchronization, with the overhead and the risk of
errors that come with it.  The model also relies on the operating system to
support kernel-level threads, which was not universal when these
architectures were first compared but is no longer a concern.

\begin{figure}[htbp]
  \centering
  % Multi-process and multithreaded web servers.  Each row of seven cells is
% the sequence of processing steps of one request.
% Usage: % Multi-process and multithreaded web servers.  Each row of seven cells is
% the sequence of processing steps of one request.
% Usage: % Multi-process and multithreaded web servers.  Each row of seven cells is
% the sequence of processing steps of one request.
% Usage: \input{figures/l06-mp-mt}   (width ~116 mm)
\begin{tikzpicture}[x=1mm, y=1mm, font=\footnotesize\sffamily,
  proc/.style={draw, rounded corners=2pt, fill=white},
  cell/.style={draw, fill=ln-box, minimum width=3.6mm, minimum height=4mm,
    inner sep=0pt},
  cache/.style={draw, fill=ln-accent!15, align=center, inner sep=0.5pt,
    font=\scriptsize\sffamily},
  lab/.style={font=\scriptsize\sffamily},
  ttl/.style={font=\footnotesize\sffamily\bfseries, anchor=west},
  >={Stealth[length=1.4mm]}]
  % (a) multi-process.  The Japanese label of the cache box is wider than the
  % English one, so the cells start further left and the box sits further in.
  \iflangja{\def\lnMPcell{12.9}\def\lnMPcache{47.9}}%
           {\def\lnMPcell{13.8}\def\lnMPcache{48}}%
  \node[ttl] at (0,31) {\iflangja{(a) マルチプロセス}{(a) multi-process}};
  \foreach \y/\n in {22/1,12/2,2/3} {
    \draw[proc] (0,\y-4) rectangle (54,\y+4);
    \node[lab] at (5.5,\y) {P\n};
    \foreach \i in {0,...,6} { \node[cell] at (\lnMPcell+4.4*\i,\y) {}; }
    \node[cache, minimum width=11mm, minimum height=6mm] at (\lnMPcache,\y)
      {\iflangja{キャッシュ}{cache}};
  }
  \node[lab, text=ln-muted, anchor=north] at (27,-2.8)
    {\iflangja{プロセスごとに別のアドレス空間}{a separate address space per process}};
  % (b) multithreaded.  Same adjustment for the wider Japanese shared-state box.
  \iflangja{\def\lnMTsh{107.6}\def\lnMTlab{65.5}\def\lnMTcell{70.1}\def\lnMTarr{98.3}}%
           {\def\lnMTsh{109.5}\def\lnMTlab{66.5}\def\lnMTcell{73.8}\def\lnMTarr{102}}%
  \node[ttl] at (62,31) {\iflangja{(b) マルチスレッド}{(b) multithreaded}};
  \draw[proc] (62,-2) rectangle (116,26);
  \node[cache, minimum width=11mm, minimum height=17mm] (sh) at (\lnMTsh,12)
    {\iflangja{共有\\状態\\（キャッシュ）}{shared\\state\\(cache)}};
  \foreach \y/\n in {19/1,12/2,5/3} {
    \node[lab] at (\lnMTlab,\y) {T\n};
    \foreach \i in {0,...,6} { \node[cell] at (\lnMTcell+4.4*\i,\y) {}; }
    \draw[<->] (\lnMTarr,\y) -- (sh.west |- 0,\y);
  }
  \node[lab, text=ln-muted, anchor=north] at (89,-2.8)
    {\iflangja{1つのアドレス空間を共有}{one shared address space}};
\end{tikzpicture}
   (width ~116 mm)
\begin{tikzpicture}[x=1mm, y=1mm, font=\footnotesize\sffamily,
  proc/.style={draw, rounded corners=2pt, fill=white},
  cell/.style={draw, fill=ln-box, minimum width=3.6mm, minimum height=4mm,
    inner sep=0pt},
  cache/.style={draw, fill=ln-accent!15, align=center, inner sep=0.5pt,
    font=\scriptsize\sffamily},
  lab/.style={font=\scriptsize\sffamily},
  ttl/.style={font=\footnotesize\sffamily\bfseries, anchor=west},
  >={Stealth[length=1.4mm]}]
  % (a) multi-process.  The Japanese label of the cache box is wider than the
  % English one, so the cells start further left and the box sits further in.
  \iflangja{\def\lnMPcell{12.9}\def\lnMPcache{47.9}}%
           {\def\lnMPcell{13.8}\def\lnMPcache{48}}%
  \node[ttl] at (0,31) {\iflangja{(a) マルチプロセス}{(a) multi-process}};
  \foreach \y/\n in {22/1,12/2,2/3} {
    \draw[proc] (0,\y-4) rectangle (54,\y+4);
    \node[lab] at (5.5,\y) {P\n};
    \foreach \i in {0,...,6} { \node[cell] at (\lnMPcell+4.4*\i,\y) {}; }
    \node[cache, minimum width=11mm, minimum height=6mm] at (\lnMPcache,\y)
      {\iflangja{キャッシュ}{cache}};
  }
  \node[lab, text=ln-muted, anchor=north] at (27,-2.8)
    {\iflangja{プロセスごとに別のアドレス空間}{a separate address space per process}};
  % (b) multithreaded.  Same adjustment for the wider Japanese shared-state box.
  \iflangja{\def\lnMTsh{107.6}\def\lnMTlab{65.5}\def\lnMTcell{70.1}\def\lnMTarr{98.3}}%
           {\def\lnMTsh{109.5}\def\lnMTlab{66.5}\def\lnMTcell{73.8}\def\lnMTarr{102}}%
  \node[ttl] at (62,31) {\iflangja{(b) マルチスレッド}{(b) multithreaded}};
  \draw[proc] (62,-2) rectangle (116,26);
  \node[cache, minimum width=11mm, minimum height=17mm] (sh) at (\lnMTsh,12)
    {\iflangja{共有\\状態\\（キャッシュ）}{shared\\state\\(cache)}};
  \foreach \y/\n in {19/1,12/2,5/3} {
    \node[lab] at (\lnMTlab,\y) {T\n};
    \foreach \i in {0,...,6} { \node[cell] at (\lnMTcell+4.4*\i,\y) {}; }
    \draw[<->] (\lnMTarr,\y) -- (sh.west |- 0,\y);
  }
  \node[lab, text=ln-muted, anchor=north] at (89,-2.8)
    {\iflangja{1つのアドレス空間を共有}{one shared address space}};
\end{tikzpicture}
   (width ~116 mm)
\begin{tikzpicture}[x=1mm, y=1mm, font=\footnotesize\sffamily,
  proc/.style={draw, rounded corners=2pt, fill=white},
  cell/.style={draw, fill=ln-box, minimum width=3.6mm, minimum height=4mm,
    inner sep=0pt},
  cache/.style={draw, fill=ln-accent!15, align=center, inner sep=0.5pt,
    font=\scriptsize\sffamily},
  lab/.style={font=\scriptsize\sffamily},
  ttl/.style={font=\footnotesize\sffamily\bfseries, anchor=west},
  >={Stealth[length=1.4mm]}]
  % (a) multi-process.  The Japanese label of the cache box is wider than the
  % English one, so the cells start further left and the box sits further in.
  \iflangja{\def\lnMPcell{12.9}\def\lnMPcache{47.9}}%
           {\def\lnMPcell{13.8}\def\lnMPcache{48}}%
  \node[ttl] at (0,31) {\iflangja{(a) マルチプロセス}{(a) multi-process}};
  \foreach \y/\n in {22/1,12/2,2/3} {
    \draw[proc] (0,\y-4) rectangle (54,\y+4);
    \node[lab] at (5.5,\y) {P\n};
    \foreach \i in {0,...,6} { \node[cell] at (\lnMPcell+4.4*\i,\y) {}; }
    \node[cache, minimum width=11mm, minimum height=6mm] at (\lnMPcache,\y)
      {\iflangja{キャッシュ}{cache}};
  }
  \node[lab, text=ln-muted, anchor=north] at (27,-2.8)
    {\iflangja{プロセスごとに別のアドレス空間}{a separate address space per process}};
  % (b) multithreaded.  Same adjustment for the wider Japanese shared-state box.
  \iflangja{\def\lnMTsh{107.6}\def\lnMTlab{65.5}\def\lnMTcell{70.1}\def\lnMTarr{98.3}}%
           {\def\lnMTsh{109.5}\def\lnMTlab{66.5}\def\lnMTcell{73.8}\def\lnMTarr{102}}%
  \node[ttl] at (62,31) {\iflangja{(b) マルチスレッド}{(b) multithreaded}};
  \draw[proc] (62,-2) rectangle (116,26);
  \node[cache, minimum width=11mm, minimum height=17mm] (sh) at (\lnMTsh,12)
    {\iflangja{共有\\状態\\（キャッシュ）}{shared\\state\\(cache)}};
  \foreach \y/\n in {19/1,12/2,5/3} {
    \node[lab] at (\lnMTlab,\y) {T\n};
    \foreach \i in {0,...,6} { \node[cell] at (\lnMTcell+4.4*\i,\y) {}; }
    \draw[<->] (\lnMTarr,\y) -- (sh.west |- 0,\y);
  }
  \node[lab, text=ln-muted, anchor=north] at (89,-2.8)
    {\iflangja{1つのアドレス空間を共有}{one shared address space}};
\end{tikzpicture}

  \caption{(a) The multi-process model: every process executes all
    processing steps (cells) of one request and keeps its own copy of state
    such as a cache.  (b) The multithreaded model: threads in one address
    space execute the steps and share the state.}
  \label{fig:l06-mp-mt}
\end{figure}

\needspace{6\baselineskip}
\section{The event-driven model}
\label{sec:l06-event}

A third architecture achieves concurrency without multiple execution
contexts.\source{P2L5, event-driven model; OSTEP ch.~33; Kerrisk
\cite{kerrisk2010} ch.~63.}

\begin{definition}[Event-driven model]
  An application that runs as a single process with a single thread of
  control in a single address space, and performs its work in response to
  events, follows the \term{event-driven model}[イベント駆動モデル].
\end{definition}

At the centre of such an application is the
\term{event dispatcher}[イベントディスパッチャ], a loop that
repeatedly waits for events and, depending on the event, invokes one or more
of the handlers registered for it (\cref{fig:l06-event-driven}).  For a web
server the events include the arrival of a connection request, the arrival
of request data on a connection, the completion of a send, which means that
the connection can take more data, and the completion of a disk read.
Pai et al.\ call a server built in this way an instance of the
\term{single-process event-driven model}[単一プロセスイベント駆動モデル]%
\index{SPED|see{single-process event-driven model}} (SPED).

\begin{figure}[htbp]
  \centering
  % Event-driven web server: one dispatcher loop invokes the handler for each
% event and keeps the state of every request explicitly.
% Usage: % Event-driven web server: one dispatcher loop invokes the handler for each
% event and keeps the state of every request explicitly.
% Usage: % Event-driven web server: one dispatcher loop invokes the handler for each
% event and keeps the state of every request explicitly.
% Usage: \input{figures/l06-event-driven}   (width ~116 mm)
\begin{tikzpicture}[x=1mm, y=1mm, font=\footnotesize\sffamily,
  box/.style={draw, rounded corners=2pt, fill=ln-box, align=center, inner sep=1pt},
  hd/.style={box, minimum width=36mm, minimum height=6mm, font=\scriptsize\sffamily},
  lab/.style={font=\scriptsize\sffamily, text=ln-muted, align=center},
  >={Stealth[length=1.6mm]}]
  % event sources
  \node[box, fill=white, minimum width=17mm, minimum height=9mm, font=\scriptsize\sffamily]
    (net) at (8.5,12) {\iflangja{ネットワーク\\（ソケット）}{network\\(sockets)}};
  \node[box, fill=white, minimum width=17mm, minimum height=9mm, font=\scriptsize\sffamily]
    (dsk) at (8.5,-4) {\iflangja{ディスク\\（ファイル）}{disk\\(files)}};
  % process
  \draw[ln-rule, rounded corners=3pt] (22,-20) rectangle (116,24);
  \node[lab, anchor=north west] at (23,23)
    {\iflangja{1 プロセス・1 スレッド・1 アドレス空間}{one process, one thread, one address space}};
  \node[box, fill=ln-accent!15, minimum width=30mm, minimum height=14mm]
    (disp) at (43,6) {\iflangja{イベントディスパッチャ}{event dispatcher}\\[0.5mm]
      \scriptsize\iflangja{イベントを待つループ}{loop: wait for events}\\
      \scriptsize(\texttt{epoll}, \texttt{select}, \texttt{poll})};
  \draw[->] (net.east) -- node[lab, above, sloped] {FD} (disp.west |- 0,9);
  \draw[->] (dsk.east) -- node[lab, below, sloped] {FD} (disp.west |- 0,3);
  \node[lab, text=ln-muted, anchor=north, align=center] at (10,-9.5)
    {\iflangja{非同期I/O\\またはヘルパ経由}{asynchronous I/O\\or a helper}};
  % handlers
  \node[hd] (h1) at (96,16) {\iflangja{接続の受け付け}{accept connection}};
  \node[hd] (h2) at (96,8)  {\iflangja{要求の読み込みと解析}{read and parse request}};
  \node[hd] (h3) at (96,0)  {\iflangja{ヘッダの送信}{send header}};
  \node[hd] (h4) at (96,-8) {\iflangja{ファイル読込・データ送信}{read file, send data}};
  \foreach \h in {h1,h2,h3,h4} { \draw[->] (disp.east) -- (\h.west); }
  \node[lab, anchor=north] at (95,-11.8)
    {\iflangja{ハンドラは最後まで実行して戻る}{handlers run to completion and return}};
  % per-request state
  \node[box, fill=white, minimum width=30mm, align=left, font=\scriptsize\sffamily,
    anchor=north] (st) at (43,-4.5)
    {C1: \iflangja{要求を待つ}{waiting for request}\\
     C2: \iflangja{送信を待つ}{waiting to send}\\
     C3: \iflangja{ファイルを待つ}{waiting for file data}};
  \node[lab, anchor=north] at (st.south) {\iflangja{要求ごとの状態}{per-request state}};
\end{tikzpicture}
   (width ~116 mm)
\begin{tikzpicture}[x=1mm, y=1mm, font=\footnotesize\sffamily,
  box/.style={draw, rounded corners=2pt, fill=ln-box, align=center, inner sep=1pt},
  hd/.style={box, minimum width=36mm, minimum height=6mm, font=\scriptsize\sffamily},
  lab/.style={font=\scriptsize\sffamily, text=ln-muted, align=center},
  >={Stealth[length=1.6mm]}]
  % event sources
  \node[box, fill=white, minimum width=17mm, minimum height=9mm, font=\scriptsize\sffamily]
    (net) at (8.5,12) {\iflangja{ネットワーク\\（ソケット）}{network\\(sockets)}};
  \node[box, fill=white, minimum width=17mm, minimum height=9mm, font=\scriptsize\sffamily]
    (dsk) at (8.5,-4) {\iflangja{ディスク\\（ファイル）}{disk\\(files)}};
  % process
  \draw[ln-rule, rounded corners=3pt] (22,-20) rectangle (116,24);
  \node[lab, anchor=north west] at (23,23)
    {\iflangja{1 プロセス・1 スレッド・1 アドレス空間}{one process, one thread, one address space}};
  \node[box, fill=ln-accent!15, minimum width=30mm, minimum height=14mm]
    (disp) at (43,6) {\iflangja{イベントディスパッチャ}{event dispatcher}\\[0.5mm]
      \scriptsize\iflangja{イベントを待つループ}{loop: wait for events}\\
      \scriptsize(\texttt{epoll}, \texttt{select}, \texttt{poll})};
  \draw[->] (net.east) -- node[lab, above, sloped] {FD} (disp.west |- 0,9);
  \draw[->] (dsk.east) -- node[lab, below, sloped] {FD} (disp.west |- 0,3);
  \node[lab, text=ln-muted, anchor=north, align=center] at (10,-9.5)
    {\iflangja{非同期I/O\\またはヘルパ経由}{asynchronous I/O\\or a helper}};
  % handlers
  \node[hd] (h1) at (96,16) {\iflangja{接続の受け付け}{accept connection}};
  \node[hd] (h2) at (96,8)  {\iflangja{要求の読み込みと解析}{read and parse request}};
  \node[hd] (h3) at (96,0)  {\iflangja{ヘッダの送信}{send header}};
  \node[hd] (h4) at (96,-8) {\iflangja{ファイル読込・データ送信}{read file, send data}};
  \foreach \h in {h1,h2,h3,h4} { \draw[->] (disp.east) -- (\h.west); }
  \node[lab, anchor=north] at (95,-11.8)
    {\iflangja{ハンドラは最後まで実行して戻る}{handlers run to completion and return}};
  % per-request state
  \node[box, fill=white, minimum width=30mm, align=left, font=\scriptsize\sffamily,
    anchor=north] (st) at (43,-4.5)
    {C1: \iflangja{要求を待つ}{waiting for request}\\
     C2: \iflangja{送信を待つ}{waiting to send}\\
     C3: \iflangja{ファイルを待つ}{waiting for file data}};
  \node[lab, anchor=north] at (st.south) {\iflangja{要求ごとの状態}{per-request state}};
\end{tikzpicture}
   (width ~116 mm)
\begin{tikzpicture}[x=1mm, y=1mm, font=\footnotesize\sffamily,
  box/.style={draw, rounded corners=2pt, fill=ln-box, align=center, inner sep=1pt},
  hd/.style={box, minimum width=36mm, minimum height=6mm, font=\scriptsize\sffamily},
  lab/.style={font=\scriptsize\sffamily, text=ln-muted, align=center},
  >={Stealth[length=1.6mm]}]
  % event sources
  \node[box, fill=white, minimum width=17mm, minimum height=9mm, font=\scriptsize\sffamily]
    (net) at (8.5,12) {\iflangja{ネットワーク\\（ソケット）}{network\\(sockets)}};
  \node[box, fill=white, minimum width=17mm, minimum height=9mm, font=\scriptsize\sffamily]
    (dsk) at (8.5,-4) {\iflangja{ディスク\\（ファイル）}{disk\\(files)}};
  % process
  \draw[ln-rule, rounded corners=3pt] (22,-20) rectangle (116,24);
  \node[lab, anchor=north west] at (23,23)
    {\iflangja{1 プロセス・1 スレッド・1 アドレス空間}{one process, one thread, one address space}};
  \node[box, fill=ln-accent!15, minimum width=30mm, minimum height=14mm]
    (disp) at (43,6) {\iflangja{イベントディスパッチャ}{event dispatcher}\\[0.5mm]
      \scriptsize\iflangja{イベントを待つループ}{loop: wait for events}\\
      \scriptsize(\texttt{epoll}, \texttt{select}, \texttt{poll})};
  \draw[->] (net.east) -- node[lab, above, sloped] {FD} (disp.west |- 0,9);
  \draw[->] (dsk.east) -- node[lab, below, sloped] {FD} (disp.west |- 0,3);
  \node[lab, text=ln-muted, anchor=north, align=center] at (10,-9.5)
    {\iflangja{非同期I/O\\またはヘルパ経由}{asynchronous I/O\\or a helper}};
  % handlers
  \node[hd] (h1) at (96,16) {\iflangja{接続の受け付け}{accept connection}};
  \node[hd] (h2) at (96,8)  {\iflangja{要求の読み込みと解析}{read and parse request}};
  \node[hd] (h3) at (96,0)  {\iflangja{ヘッダの送信}{send header}};
  \node[hd] (h4) at (96,-8) {\iflangja{ファイル読込・データ送信}{read file, send data}};
  \foreach \h in {h1,h2,h3,h4} { \draw[->] (disp.east) -- (\h.west); }
  \node[lab, anchor=north] at (95,-11.8)
    {\iflangja{ハンドラは最後まで実行して戻る}{handlers run to completion and return}};
  % per-request state
  \node[box, fill=white, minimum width=30mm, align=left, font=\scriptsize\sffamily,
    anchor=north] (st) at (43,-4.5)
    {C1: \iflangja{要求を待つ}{waiting for request}\\
     C2: \iflangja{送信を待つ}{waiting to send}\\
     C3: \iflangja{ファイルを待つ}{waiting for file data}};
  \node[lab, anchor=north] at (st.south) {\iflangja{要求ごとの状態}{per-request state}};
\end{tikzpicture}

  \caption{An event-driven web server.  The dispatcher waits for events on
    the file descriptors (FDs) of sockets and files and invokes the handler
    for each event.  The state of every request in progress is kept in an
    explicit record.}
  \label{fig:l06-event-driven}
\end{figure}

An \term{event handler}[イベントハンドラ] is the code executed in response
to one event; invoking it is simply a call to that code.  A handler runs to
completion.  If it would have to block, it does not wait: it initiates the
blocking operation, records how far the request has progressed, and returns
control to the dispatcher, and the completion of the operation later
arrives as another event.  Dispatcher and handlers together form a state
machine.  The information that a thread would keep implicitly, in its
program counter and on its stack, is kept explicitly in a record for each
request.\margin{OSTEP calls this explicit bookkeeping \emph{manual stack
management}.}

\subsection{Concurrent execution}

Although it has a single thread, an event-driven server processes many
requests concurrently.  Every request in progress is at some step of its
processing.  The dispatcher executes one step of one request, then turns to
the next event, which may belong to a different request, so that the
processing of the requests is interleaved.  \Cref{tab:l06-trace} traces
three connections; in step~7, for instance, one thread has three requests
in progress, each at a different step.

\begin{table}[htbp]
  \centering
  \caption{Trace of an event-driven server with three connections.  States:
    R waiting for the request, F waiting for file data, S waiting until
    the connection can take more data, --- not yet connected.}
  \label{tab:l06-trace}
  \small
  \begin{tabular}{@{}c>{\raggedright\arraybackslash}p{0.25\linewidth}>{\raggedright\arraybackslash}p{0.33\linewidth}ccc@{}}
    \toprule
    Step & Event & Handler action & C1 & C2 & C3 \\
    \midrule
    1 & connection request & accept C1 & R & --- & --- \\
    2 & connection request & accept C2 & R & R & --- \\
    3 & request on C1 & parse, find file, start read & F & R & --- \\
    4 & connection request & accept C3 & F & R & R \\
    5 & request on C2 & parse, find file, start read & F & F & R \\
    6 & file data for C1 & send header and data & S & F & R \\
    7 & request on C3 & parse, find file, start read & S & F & F \\
    8 & C1 can take data & send the rest, close & done & F & F \\
    9 & file data for C2 & send header and data & done & S & F \\
    \bottomrule
  \end{tabular}
\end{table}

\subsection{Why the model works}

Lesson~3 showed that on a single CPU, switching to another thread while a
thread waits for I/O pays off when the idle time exceeds the cost of the two
context switches, $t_{\mathrm{idle}} > 2\, t_{\mathrm{ctx}}$.  When there is
no idle time, context switches are pure overhead: they consume CPU cycles
that could have processed requests.  The event-driven model obtains the
benefit without the cost.  It processes a request until the request has to
wait, then processes another request, and no context switch is involved.
On a machine with several CPUs, the model is applied once per CPU: several
event-driven processes, each with its own dispatcher, share the incoming
connections.

\subsection{Implementation}

Both kinds of event sources of a web server, sockets for the network and
files for the disk, are accessed through file descriptors, and an event is
the readiness of a file descriptor for input or output.  To find out which
descriptors are ready, the dispatcher uses one of three system calls.
\texttt{select} and \texttt{poll} receive the whole set of descriptors of
interest on every call; the kernel checks all of them, and the application
scans the result, although typically only a few are ready.  Their cost
therefore grows with the number of open connections.  \texttt{epoll},
available on Linux, eliminates most of this overhead: descriptors are
registered once, and each call returns only those that are ready.

\begin{example}
  The skeleton of a dispatcher built on \texttt{epoll}, with the state of
  each request in a \texttt{struct conn}:
  \begin{lstlisting}[language=C, numbers=none]
struct conn {                   /* one per request */
    int fd;
    void (*handler)(struct conn *);   /* next step */
    /* buffers, file, offset, ... */
};

int ep = epoll_create1(0);
/* handlers register descriptors with epoll_ctl */
for (;;) {                                /* dispatcher */
    struct epoll_event ev[64];
    int n = epoll_wait(ep, ev, 64, -1);   /* wait */
    for (int i = 0; i < n; i++) {
        struct conn *c = ev[i].data.ptr;
        c->handler(c);           /* runs to completion */
    }
}
\end{lstlisting}
  The handler of the listening socket accepts a connection, allocates a
  \texttt{struct conn} whose handler reads the request, and registers the
  new descriptor.  Every handler performs one step and sets
  \texttt{handler} to the step that follows.
\end{example}

The benefits of the model follow from its single address space and single
flow of control.  A request in progress costs only its state record, not a
thread with its stack, so the memory requirement is small.  There are no
context switches between execution contexts.  And since handlers never run
concurrently, no synchronization is needed.

\begin{keypoint}
  An event-driven server interleaves the processing of many requests in one
  thread: a dispatcher waits for events on file descriptors and runs a
  handler for each.  It saves memory, context switches and synchronization,
  provided that no handler blocks.
\end{keypoint}

\needspace{6\baselineskip}
\section{Blocking operations and helpers}
\label{sec:l06-helpers}

The event-driven model has one serious problem: a handler that blocks
blocks the entire process.\source{P2L5, asynchronous I/O, helpers;
OSTEP ch.~33; Pai et al.}  While the single thread waits for the
disk, the dispatcher does not run, and no other request makes progress, not
even a request whose data are ready to be sent.  Readiness notification
does not help with disk reads: \texttt{select} and \texttt{poll} always
report a regular file as ready, even when reading it must wait for the
disk, and \texttt{epoll} does not accept regular files at all.

\begin{definition}[Asynchronous I/O]
  With \term{asynchronous I/O}[非同期I/O] a system call that starts an I/O
  operation returns before the operation has completed, and the calling
  thread continues to execute.  The kernel obtains from the call all the
  information it needs to carry out the operation, and either learns where
  to deliver the results or tells the caller where to retrieve them later.
\end{definition}

Asynchronous I/O requires support from the kernel, for example kernel-level
threads that carry out the blocking part of the operation, or from the
device, for example a device that transfers the data by direct memory access
(DMA, Lesson~11) and signals completion with an interrupt.  It fits the
event-driven model well: the completion of an asynchronous operation is one
more event for the dispatcher.\margin{POSIX defines \texttt{aio\_read} and
\texttt{aio\_write}; glibc implements them with threads.  Linux has also
offered \texttt{io\_uring} since version 5.1.}

Asynchronous calls are not always available, however, and where they exist
they may not be supported for every kind of device.  An event-driven server
then delegates its blocking operations.

\begin{definition}[Helper]
  A \term{helper}[ヘルパ] is a thread or process that an event-driven server
  uses only to execute blocking operations on behalf of its dispatcher.
\end{definition}

The dispatcher passes an operation to a helper through a pipe or a socket,
and the helper reports completion through the same channel.  Because pipes
and sockets are file descriptors, the dispatcher waits for these reports
with \texttt{select} or \texttt{poll} like for any other event.  The helper
blocks; the dispatcher, and with it all other requests, does not.  An
event-driven process combined with helper processes forms the
\term{asymmetric multi-process event-driven model}[非対称マルチプロセス\allowbreak イベント駆動モデル]%
\index{AMPED|see{asymmetric multi-process event-driven model}} (AMPED).
The model is \emph{asymmetric} because the helpers perform only particular
operations, unlike the equivalent processes or threads of the MP and MT
models.  Pai et al.\ used helper processes because, at the time, not every
operating system supported kernel-level threads.  With helper threads
instead of processes, the server follows the
\term{asymmetric multithreaded event-driven model}[非対称マルチスレッド\allowbreak イベント駆動モデル]%
\index{AMTED|see{asymmetric multithreaded event-driven model}} (AMTED).

Helpers resolve the portability problem of the basic event-driven model:
the server does not depend on asynchronous I/O support in the kernel.  They
also keep the memory footprint smaller than that of a worker for every
request, since helpers are needed only for the operations that are blocked
at the same time, not for every connection.  On the other hand, the
approach applies only to certain classes of applications, and on a
multiprocessor it complicates the routing of events: each event, including
the completion reported by a helper, must reach the dispatcher that handles
the corresponding request.

\begin{example}
  A server has 2\,000 open connections, of which at most 25 wait for the
  disk at any time.  In the multithreaded model each connection needs a
  thread; with \SI{64}{\kilo\byte} per thread for its stack and data
  structures, the threads use \SI{128}{\mega\byte}.  An AMPED server keeps a
  state record of, say, \SI{2}{\kilo\byte} per connection
  (\SI{4}{\mega\byte}) and 25 helpers of \SI{64}{\kilo\byte} each
  (\SI{1.6}{\mega\byte}), about \SI{5.6}{\mega\byte} in total.  The
  difference can be used to cache file data, which in turn reduces the
  number of requests that must wait for the disk.
\end{example}

\Cref{tab:l06-models} summarizes the four architectures.

\begin{table}[htbp]
  \centering
  \caption{Architectures of a concurrent web server.}
  \label{tab:l06-models}
  \small
  \begin{tabular}{@{}>{\raggedright\arraybackslash}p{0.19\linewidth}>{\raggedright\arraybackslash}p{0.165\linewidth}>{\raggedright\arraybackslash}p{0.165\linewidth}>{\raggedright\arraybackslash}p{0.165\linewidth}>{\raggedright\arraybackslash}p{0.165\linewidth}@{}}
    \toprule
    & MP & MT & SPED & AMPED \\
    \midrule
    Execution contexts & a process per concurrent request & a thread per
      concurrent request & one & one, plus helpers \\
    Memory per request & address space & thread stack and structures &
      state record & state record \\
    Context switches & between processes & between threads & none & only
      for blocking operations \\
    Synchronization & through IPC & locks on shared state & none & none \\
    Blocking disk I/O & blocks one process & blocks one thread & blocks the
      server & blocks one helper \\
    \bottomrule
  \end{tabular}
\end{table}

\needspace{6\baselineskip}
\section{The Flash web server}
\label{sec:l06-flash}

Pai et al.\ built an event-driven web server based on the AMPED
model.\source{P2L5, Flash; Pai et al.; Kerrisk ch.~49--50.}  In their server,
\term{Flash}[Flash], helpers perform the disk operations that may block,
and they communicate with the dispatcher through pipes.  Flash accesses file data through memory-mapped
files (Lesson~9), so a page of a file is read from disk when it is first
accessed.

Before a handler sends file data, the dispatcher uses the system call
\texttt{mincore} to test whether the pages of the data are in memory
(\cref{fig:l06-flash}).  If they are, a handler sends the data at once, and
this cannot block.  If they are not, the dispatcher passes the request to a
helper through a pipe.  The helper accesses the pages, which brings them
into memory and may block, and then reports completion through the pipe.
The dispatcher receives the report as an event and invokes the handler,
which now finds the data in memory.  Most requests to a web server are for
popular files whose data are already in memory; for these requests Flash
avoids the helpers and the cost of communicating with them, which can be a
large saving.

\begin{figure}[htbp]
  \centering
  % Flash (AMPED): the dispatcher tests with mincore whether the file data are
% in memory; if not, a helper process reads them and reports completion
% through a pipe.
% Usage: % Flash (AMPED): the dispatcher tests with mincore whether the file data are
% in memory; if not, a helper process reads them and reports completion
% through a pipe.
% Usage: % Flash (AMPED): the dispatcher tests with mincore whether the file data are
% in memory; if not, a helper process reads them and reports completion
% through a pipe.
% Usage: \input{figures/l06-flash}   (width ~116 mm)
\begin{tikzpicture}[x=1mm, y=1mm, font=\footnotesize\sffamily,
  box/.style={draw, rounded corners=2pt, fill=ln-box, align=center, inner sep=1pt},
  lab/.style={font=\scriptsize\sffamily, align=center},
  num/.style={draw, circle, fill=white, inner sep=0pt, minimum size=4mm,
    font=\scriptsize\sffamily\bfseries, text=ln-accent},
  >={Stealth[length=1.6mm]}]
  % main process
  \draw[ln-rule, rounded corners=3pt] (0,-15) rectangle (64,21);
  \node[lab, text=ln-muted, anchor=north west] at (1,20)
    {\iflangja{Flash のプロセス（ディスパッチャとハンドラ）}{Flash process (dispatcher and handlers)}};
  \node[box, fill=ln-accent!15, minimum width=20mm, minimum height=10mm]
    (disp) at (12,6) {\iflangja{ディスパッチャ}{dispatcher}};
  \node[box, fill=white, minimum width=24mm, minimum height=10mm, font=\scriptsize\sffamily]
    (chk) at (48,6) {\iflangja{ページはメモリ上か}{pages in memory?}\\(\texttt{mincore})};
  \node[box, minimum width=24mm, minimum height=8mm, font=\scriptsize\sffamily]
    (hnd) at (48,-8) {\iflangja{ハンドラ：}{handler:}\\\iflangja{データを送信}{send data}};
  \draw[->] (disp) -- node[num, above=0.6mm] {1} (chk);
  \draw[->] (chk) -- node[lab, right] {\iflangja{はい}{yes}} node[num, left=0.6mm] {2} (hnd);
  % helpers
  \foreach \y/\n in {22/1,14/2,6/3} {
    \node[box, minimum width=16mm, minimum height=6mm, font=\scriptsize\sffamily]
      (h\n) at (86,\y) {\iflangja{ヘルパ}{helper} \n};
  }
  \draw[->] (chk) -- node[lab, above, pos=0.62] {\iflangja{いいえ：\\パイプ}{no: pipe}}
    node[num, below=0.6mm, pos=0.62] {3} (h3);
  % disk
  \node[cylinder, shape border rotate=90, draw, aspect=0.25, fill=ln-box,
    minimum width=11mm, minimum height=13mm, font=\scriptsize\sffamily]
    (disk) at (109,6) {\iflangja{ディスク}{disk}};
  \draw[<->] (h3.east) -- node[lab, above] {\texttt{mmap}} (disk.west);
  % completion
  \draw[->, densely dashed, ln-caution] (h3.south) -- (86,-20) -- (12,-20)
    -- node[num, left=0.6mm, pos=0.5] {4} (disp.south);
  \node[lab, text=ln-caution, anchor=south] at (48,-19.6)
    {\iflangja{完了をパイプで通知（イベント）}{completion reported through a pipe (an event)}};
  \draw[->, densely dashed, ln-caution] (disp.south east) -- (hnd.west);
\end{tikzpicture}
   (width ~116 mm)
\begin{tikzpicture}[x=1mm, y=1mm, font=\footnotesize\sffamily,
  box/.style={draw, rounded corners=2pt, fill=ln-box, align=center, inner sep=1pt},
  lab/.style={font=\scriptsize\sffamily, align=center},
  num/.style={draw, circle, fill=white, inner sep=0pt, minimum size=4mm,
    font=\scriptsize\sffamily\bfseries, text=ln-accent},
  >={Stealth[length=1.6mm]}]
  % main process
  \draw[ln-rule, rounded corners=3pt] (0,-15) rectangle (64,21);
  \node[lab, text=ln-muted, anchor=north west] at (1,20)
    {\iflangja{Flash のプロセス（ディスパッチャとハンドラ）}{Flash process (dispatcher and handlers)}};
  \node[box, fill=ln-accent!15, minimum width=20mm, minimum height=10mm]
    (disp) at (12,6) {\iflangja{ディスパッチャ}{dispatcher}};
  \node[box, fill=white, minimum width=24mm, minimum height=10mm, font=\scriptsize\sffamily]
    (chk) at (48,6) {\iflangja{ページはメモリ上か}{pages in memory?}\\(\texttt{mincore})};
  \node[box, minimum width=24mm, minimum height=8mm, font=\scriptsize\sffamily]
    (hnd) at (48,-8) {\iflangja{ハンドラ：}{handler:}\\\iflangja{データを送信}{send data}};
  \draw[->] (disp) -- node[num, above=0.6mm] {1} (chk);
  \draw[->] (chk) -- node[lab, right] {\iflangja{はい}{yes}} node[num, left=0.6mm] {2} (hnd);
  % helpers
  \foreach \y/\n in {22/1,14/2,6/3} {
    \node[box, minimum width=16mm, minimum height=6mm, font=\scriptsize\sffamily]
      (h\n) at (86,\y) {\iflangja{ヘルパ}{helper} \n};
  }
  \draw[->] (chk) -- node[lab, above, pos=0.62] {\iflangja{いいえ：\\パイプ}{no: pipe}}
    node[num, below=0.6mm, pos=0.62] {3} (h3);
  % disk
  \node[cylinder, shape border rotate=90, draw, aspect=0.25, fill=ln-box,
    minimum width=11mm, minimum height=13mm, font=\scriptsize\sffamily]
    (disk) at (109,6) {\iflangja{ディスク}{disk}};
  \draw[<->] (h3.east) -- node[lab, above] {\texttt{mmap}} (disk.west);
  % completion
  \draw[->, densely dashed, ln-caution] (h3.south) -- (86,-20) -- (12,-20)
    -- node[num, left=0.6mm, pos=0.5] {4} (disp.south);
  \node[lab, text=ln-caution, anchor=south] at (48,-19.6)
    {\iflangja{完了をパイプで通知（イベント）}{completion reported through a pipe (an event)}};
  \draw[->, densely dashed, ln-caution] (disp.south east) -- (hnd.west);
\end{tikzpicture}
   (width ~116 mm)
\begin{tikzpicture}[x=1mm, y=1mm, font=\footnotesize\sffamily,
  box/.style={draw, rounded corners=2pt, fill=ln-box, align=center, inner sep=1pt},
  lab/.style={font=\scriptsize\sffamily, align=center},
  num/.style={draw, circle, fill=white, inner sep=0pt, minimum size=4mm,
    font=\scriptsize\sffamily\bfseries, text=ln-accent},
  >={Stealth[length=1.6mm]}]
  % main process
  \draw[ln-rule, rounded corners=3pt] (0,-15) rectangle (64,21);
  \node[lab, text=ln-muted, anchor=north west] at (1,20)
    {\iflangja{Flash のプロセス（ディスパッチャとハンドラ）}{Flash process (dispatcher and handlers)}};
  \node[box, fill=ln-accent!15, minimum width=20mm, minimum height=10mm]
    (disp) at (12,6) {\iflangja{ディスパッチャ}{dispatcher}};
  \node[box, fill=white, minimum width=24mm, minimum height=10mm, font=\scriptsize\sffamily]
    (chk) at (48,6) {\iflangja{ページはメモリ上か}{pages in memory?}\\(\texttt{mincore})};
  \node[box, minimum width=24mm, minimum height=8mm, font=\scriptsize\sffamily]
    (hnd) at (48,-8) {\iflangja{ハンドラ：}{handler:}\\\iflangja{データを送信}{send data}};
  \draw[->] (disp) -- node[num, above=0.6mm] {1} (chk);
  \draw[->] (chk) -- node[lab, right] {\iflangja{はい}{yes}} node[num, left=0.6mm] {2} (hnd);
  % helpers
  \foreach \y/\n in {22/1,14/2,6/3} {
    \node[box, minimum width=16mm, minimum height=6mm, font=\scriptsize\sffamily]
      (h\n) at (86,\y) {\iflangja{ヘルパ}{helper} \n};
  }
  \draw[->] (chk) -- node[lab, above, pos=0.62] {\iflangja{いいえ：\\パイプ}{no: pipe}}
    node[num, below=0.6mm, pos=0.62] {3} (h3);
  % disk
  \node[cylinder, shape border rotate=90, draw, aspect=0.25, fill=ln-box,
    minimum width=11mm, minimum height=13mm, font=\scriptsize\sffamily]
    (disk) at (109,6) {\iflangja{ディスク}{disk}};
  \draw[<->] (h3.east) -- node[lab, above] {\texttt{mmap}} (disk.west);
  % completion
  \draw[->, densely dashed, ln-caution] (h3.south) -- (86,-20) -- (12,-20)
    -- node[num, left=0.6mm, pos=0.5] {4} (disp.south);
  \node[lab, text=ln-caution, anchor=south] at (48,-19.6)
    {\iflangja{完了をパイプで通知（イベント）}{completion reported through a pipe (an event)}};
  \draw[->, densely dashed, ln-caution] (disp.south east) -- (hnd.west);
\end{tikzpicture}

  \caption{Serving file data in Flash.  (1)~The dispatcher tests with
    \texttt{mincore} whether the data are in memory.  (2)~If they are, a
    handler sends them directly.  (3)~Otherwise a helper is asked through a
    pipe to access the pages, reading them from disk.  (4)~Its completion
    report is an event, after which the handler sends the data.}
  \label{fig:l06-flash}
\end{figure}

\begin{example}
  The test in C, using the headers \texttt{sys/mman.h},
  \texttt{unistd.h} and \texttt{stdbool.h}, for file data mapped at a
  page-aligned address:
  \begin{lstlisting}[language=C, numbers=none]
static bool in_memory(void *addr, size_t len)
{
    size_t pg = (size_t)sysconf(_SC_PAGESIZE);
    size_t n = (len + pg - 1) / pg;
    unsigned char vec[n];            /* 1 byte per page */
    if (mincore(addr, len, vec) != 0)
        return false;
    for (size_t i = 0; i < n; i++)
        if (!(vec[i] & 1))           /* page not resident */
            return false;
    return true;
}

/* in the handler that is about to send file data */
if (in_memory(c->map, c->len))
    send_data(c);                    /* cannot block */
else
    ask_helper(c);                   /* through a pipe */
\end{lstlisting}
\end{example}

\subsection{Additional optimizations}

Flash combines the architecture with optimizations that reduce the work per
request.
\begin{description}
  \item[Application-level caching] Flash caches both data and the results
    of computations: the translation of the path name in a request to a
    file, the response headers, which depend only on the file, and the
    mapped files themselves.
  \item[Alignment for DMA] Flash places the data it sends in memory so that
    a network interface capable of DMA can transfer them directly, without
    copying them into kernel buffers.
  \item[Gather I/O] The response header and the file data are sent with a
    single vector I/O operation, \texttt{writev}, which passes several
    memory regions at once; a network interface with scatter-gather DMA can
    transfer the regions without first assembling them in one buffer.
\end{description}
These optimizations were novel in 1999 and are common in web servers today.
They are independent of the server architecture, a fact that matters for
the evaluation in \cref{sec:l06-methodology}.

\needspace{6\baselineskip}
\section{The Apache web server}
\label{sec:l06-apache}

The open-source \term{Apache}[Apache] HTTP server shows how a production
server combines the models.\source{P2L5, Apache web server.}  Its
\emph{core} provides the basic server skeleton: it accepts connections and
manages the processes and threads.  All other functionality is implemented
by \emph{modules}, each of which registers handlers at \emph{hooks}, the
stages of request processing (\cref{fig:l06-apache}a).  As a request
progresses, the core invokes the handlers registered for each stage, a flow
of control that resembles the event-driven model; unlike event handlers,
however, these handlers run on the thread assigned to the request and may
block.

\begin{figure}[htbp]
  \centering
  % Apache: (a) the core invokes modules registered for the stages of request
% processing; (b) several processes, each a boss with a dynamic worker pool.
% Usage: % Apache: (a) the core invokes modules registered for the stages of request
% processing; (b) several processes, each a boss with a dynamic worker pool.
% Usage: % Apache: (a) the core invokes modules registered for the stages of request
% processing; (b) several processes, each a boss with a dynamic worker pool.
% Usage: \input{figures/l06-apache}   (width ~116 mm)
\begin{tikzpicture}[x=1mm, y=1mm, font=\footnotesize\sffamily,
  box/.style={draw, rounded corners=2pt, fill=ln-box, align=center, inner sep=1pt},
  mod/.style={box, fill=white, minimum width=11mm, minimum height=8mm,
    font=\scriptsize\sffamily},
  th/.style={draw, circle, minimum size=5mm, inner sep=0pt,
    font=\scriptsize\sffamily},
  lab/.style={font=\scriptsize\sffamily, text=ln-muted, align=center},
  ttl/.style={font=\footnotesize\sffamily\bfseries, anchor=west},
  >={Stealth[length=1.4mm]}]
  % (a) core and modules
  \node[ttl] at (0,31) {\iflangja{(a) 要求の処理}{(a) request processing}};
  \node[box, fill=ln-accent!15, minimum width=48mm, minimum height=7mm]
    (core) at (24,20) {\iflangja{コア：処理の各段階}{core: stages of processing}};
  \foreach \x/\t in {6/{\iflangja{認証}{authen-\\tication}},
                     18/{\iflangja{ファイル\\提供}{file\\serving}},
                     30/{CGI},
                     42/{\iflangja{ログ}{logging}}} {
    \node[mod] (m\x) at (\x,5) {\t};
    \draw[<->] (\x,16.5) -- (m\x.north);
  }
  \node[lab] at (24,-3) {\iflangja{モジュール（フックに登録）}{modules, registered at hooks}};
  % (b) processes
  \node[ttl] at (56,31) {\iflangja{(b) プロセスとスレッド}{(b) processes and threads}};
  \foreach \px/\n in {56/1,84/2} {
    \draw[rounded corners=2pt] (\px,0) rectangle (\px+25,26);
    \node[th, fill=ln-accent!25] at (\px+12.5,20.5) {B};
    \foreach \i in {0,1,2} {
      \node[th, fill=ln-box] at (\px+5.5+7*\i,12.5) {W};
    }
    \node[th, fill=ln-box] at (\px+5.5,6) {W};
    \node[th, densely dashed] at (\px+12.5,6) {};
    \node[th, densely dashed] at (\px+19.5,6) {};
    \foreach \i in {0,1,2} { \draw[->] (\px+12.5,18) -- (\px+5.5+7*\i,15.2); }
    \node[font=\scriptsize\sffamily, text=ln-muted] at (\px+12.5,1.6)
      {\iflangja{プロセス}{process} \n};
  }
  \node at (113,13) {\dots};
  \node[lab] at (86,-3) {\iflangja{プロセス数もスレッド数も負荷に応じて増減}{processes and thread pools grow and shrink with load}};
\end{tikzpicture}
   (width ~116 mm)
\begin{tikzpicture}[x=1mm, y=1mm, font=\footnotesize\sffamily,
  box/.style={draw, rounded corners=2pt, fill=ln-box, align=center, inner sep=1pt},
  mod/.style={box, fill=white, minimum width=11mm, minimum height=8mm,
    font=\scriptsize\sffamily},
  th/.style={draw, circle, minimum size=5mm, inner sep=0pt,
    font=\scriptsize\sffamily},
  lab/.style={font=\scriptsize\sffamily, text=ln-muted, align=center},
  ttl/.style={font=\footnotesize\sffamily\bfseries, anchor=west},
  >={Stealth[length=1.4mm]}]
  % (a) core and modules
  \node[ttl] at (0,31) {\iflangja{(a) 要求の処理}{(a) request processing}};
  \node[box, fill=ln-accent!15, minimum width=48mm, minimum height=7mm]
    (core) at (24,20) {\iflangja{コア：処理の各段階}{core: stages of processing}};
  \foreach \x/\t in {6/{\iflangja{認証}{authen-\\tication}},
                     18/{\iflangja{ファイル\\提供}{file\\serving}},
                     30/{CGI},
                     42/{\iflangja{ログ}{logging}}} {
    \node[mod] (m\x) at (\x,5) {\t};
    \draw[<->] (\x,16.5) -- (m\x.north);
  }
  \node[lab] at (24,-3) {\iflangja{モジュール（フックに登録）}{modules, registered at hooks}};
  % (b) processes
  \node[ttl] at (56,31) {\iflangja{(b) プロセスとスレッド}{(b) processes and threads}};
  \foreach \px/\n in {56/1,84/2} {
    \draw[rounded corners=2pt] (\px,0) rectangle (\px+25,26);
    \node[th, fill=ln-accent!25] at (\px+12.5,20.5) {B};
    \foreach \i in {0,1,2} {
      \node[th, fill=ln-box] at (\px+5.5+7*\i,12.5) {W};
    }
    \node[th, fill=ln-box] at (\px+5.5,6) {W};
    \node[th, densely dashed] at (\px+12.5,6) {};
    \node[th, densely dashed] at (\px+19.5,6) {};
    \foreach \i in {0,1,2} { \draw[->] (\px+12.5,18) -- (\px+5.5+7*\i,15.2); }
    \node[font=\scriptsize\sffamily, text=ln-muted] at (\px+12.5,1.6)
      {\iflangja{プロセス}{process} \n};
  }
  \node at (113,13) {\dots};
  \node[lab] at (86,-3) {\iflangja{プロセス数もスレッド数も負荷に応じて増減}{processes and thread pools grow and shrink with load}};
\end{tikzpicture}
   (width ~116 mm)
\begin{tikzpicture}[x=1mm, y=1mm, font=\footnotesize\sffamily,
  box/.style={draw, rounded corners=2pt, fill=ln-box, align=center, inner sep=1pt},
  mod/.style={box, fill=white, minimum width=11mm, minimum height=8mm,
    font=\scriptsize\sffamily},
  th/.style={draw, circle, minimum size=5mm, inner sep=0pt,
    font=\scriptsize\sffamily},
  lab/.style={font=\scriptsize\sffamily, text=ln-muted, align=center},
  ttl/.style={font=\footnotesize\sffamily\bfseries, anchor=west},
  >={Stealth[length=1.4mm]}]
  % (a) core and modules
  \node[ttl] at (0,31) {\iflangja{(a) 要求の処理}{(a) request processing}};
  \node[box, fill=ln-accent!15, minimum width=48mm, minimum height=7mm]
    (core) at (24,20) {\iflangja{コア：処理の各段階}{core: stages of processing}};
  \foreach \x/\t in {6/{\iflangja{認証}{authen-\\tication}},
                     18/{\iflangja{ファイル\\提供}{file\\serving}},
                     30/{CGI},
                     42/{\iflangja{ログ}{logging}}} {
    \node[mod] (m\x) at (\x,5) {\t};
    \draw[<->] (\x,16.5) -- (m\x.north);
  }
  \node[lab] at (24,-3) {\iflangja{モジュール（フックに登録）}{modules, registered at hooks}};
  % (b) processes
  \node[ttl] at (56,31) {\iflangja{(b) プロセスとスレッド}{(b) processes and threads}};
  \foreach \px/\n in {56/1,84/2} {
    \draw[rounded corners=2pt] (\px,0) rectangle (\px+25,26);
    \node[th, fill=ln-accent!25] at (\px+12.5,20.5) {B};
    \foreach \i in {0,1,2} {
      \node[th, fill=ln-box] at (\px+5.5+7*\i,12.5) {W};
    }
    \node[th, fill=ln-box] at (\px+5.5,6) {W};
    \node[th, densely dashed] at (\px+12.5,6) {};
    \node[th, densely dashed] at (\px+19.5,6) {};
    \foreach \i in {0,1,2} { \draw[->] (\px+12.5,18) -- (\px+5.5+7*\i,15.2); }
    \node[font=\scriptsize\sffamily, text=ln-muted] at (\px+12.5,1.6)
      {\iflangja{プロセス}{process} \n};
  }
  \node at (113,13) {\dots};
  \node[lab] at (86,-3) {\iflangja{プロセス数もスレッド数も負荷に応じて増減}{processes and thread pools grow and shrink with load}};
\end{tikzpicture}

  \caption{The Apache web server.  (a) Modules register handlers at the
    stages of request processing that the core provides.  (b) Several
    processes, each a boss (B) with a pool of workers (W); dashed circles are
    workers created as load increases.}
  \label{fig:l06-apache}
\end{figure}

The execution contexts combine the multi-process and multithreaded models
(\cref{fig:l06-apache}b).  Apache runs several processes, and each process
is a boss-workers server with a dynamic thread pool.  Both the number of
threads in each pool and the number of processes are adjusted to the load,
for instance by the number of idle workers or of pending
connections.

\begin{note}
  In Apache~2 the structure of the execution contexts is selected by a
  \emph{multi-processing module}: \texttt{prefork} uses processes only,
  \texttt{worker} is the hybrid described here, and \texttt{event} adds an
  event loop that watches idle connections.
\end{note}

\needspace{6\baselineskip}
\section{Experimental methodology}
\label{sec:l06-methodology}

An experimental comparison must answer three questions: which systems are
compared, with which inputs, and by which metrics.\source{P2L5, experimental
methodology; Pai et al.}  The evaluation of Flash illustrates each.

\subsection{Comparison points}

Flash was compared with servers of every architecture discussed in this
lesson (\cref{tab:l06-comparison-points}).  The MP, MT and SPED servers
were built from the code of Flash and differ from it only in their
architecture; they share its optimizations.  Differences among them can
therefore be attributed to the architecture.  Two independent servers show
where Flash stands among the servers in use at the time.

\begin{table}[htbp]
  \centering
  \caption{Servers compared in the evaluation of Flash.}
  \label{tab:l06-comparison-points}
  \small
  \begin{tabular}{@{}l>{\raggedright\arraybackslash}p{0.72\linewidth}@{}}
    \toprule
    Server & Architecture \\
    \midrule
    MP     & multi-process: each process has a single thread and handles one
             request at a time \\
    MT     & multithreaded, boss-workers \\
    SPED   & single-process event-driven, without helpers \\
    Zeus   & a commercial SPED server, run with two processes \\
    Apache & version 1.3.1, multi-process \\
    Flash  & AMPED \\
    \bottomrule
  \end{tabular}
\end{table}

\subsection{Workloads}

\begin{definition}[Workload]
  The \term{workload}[ワークロード] of an experiment is the set of inputs
  presented to the system under test; for a web server, the sequence of
  requests with their arrival times and the files they request.
\end{definition}

A workload should be realistic, reproducing for instance the distribution
of web page requests over time that real servers see, and it should be
controlled and reproducible, so that every server is measured with the same
input.  A \term{trace}[トレース], a record of the requests observed by a
real system that the experiment replays, meets both requirements.  A
\term{synthetic workload}[合成ワークロード] is generated by a program from
chosen parameters; it is less realistic, but it lets each parameter be
varied systematically.  Flash was evaluated with two traces collected at
web servers of Rice University: the \emph{CS trace} of the computer science
department, with a large set of files, and the \emph{Owlnet trace} of a
server hosting the personal pages of students, with a smaller set.  A
synthetic workload repeatedly requested a single file whose size was varied.

\subsection{Metrics}

A web server is measured by the data it delivers and by the clients it
serves, in two metrics:
\begin{equation}
  \begin{aligned}
    \text{bandwidth} &= \frac{\text{total bytes transferred}}{\text{total time}} ,\\[1mm]
    \text{connection rate} &= \frac{\text{total client connections}}{\text{total time}} .
  \end{aligned}
  \label{eq:l06-bandwidth}
\end{equation}
Both were evaluated as functions of the file size.  A larger file lets the
server amortize the fixed cost of a connection over more bytes, which raises
the bandwidth, but it also means more work per connection, which lowers the
connection rate.  A simple model makes this precise.  If serving a file of
size $S$ costs a fixed time $c$ per connection plus $b$ per byte, a server
that handles one connection at a time achieves
\begin{equation}
  R(S) = \frac{1}{c + b\,S} ,
  \qquad
  B(S) = S \cdot R(S) = \frac{S}{c + b\,S} \;<\; \frac{1}{b} ,
  \label{eq:l06-amortize}
\end{equation}
where $R$ is the connection rate and $B$ the bandwidth.

\begin{example}
  With $c = \SI{1}{\milli\second}$ and $b = \SI{0.02}{\milli\second}$ per
  kilobyte, a file of \SI{10}{\kilo\byte} takes \SI{1.2}{\milli\second}:
  $R \approx 833$ connections per second and
  $B \approx \SI[per-mode=symbol]{8.3}{\mega\byte\per\second}$.  A file of
  \SI{100}{\kilo\byte} takes \SI{3}{\milli\second}: $R \approx 333$
  connections per second, but $B \approx \SI[per-mode=symbol]{33.3}{\mega\byte\per\second}$.
  As the file grows, $B$ approaches $1/b = \SI[per-mode=symbol]{50}{\mega\byte\per\second}$
  while $R$ keeps falling.
\end{example}

\needspace{6\baselineskip}
\section{Experimental results}
\label{sec:l06-results}

The results of the evaluation depend above all on whether the requested
data are in memory.\source{P2L5, experimental results, summary of
performance results; Pai et al.}
\Cref{tab:l06-results} summarizes the findings described below.

\begin{table}[htbp]
  \centering
  \caption{Summary of the performance results of the Flash evaluation.}
  \label{tab:l06-results}
  \small
  \begin{tabular}{@{}>{\raggedright\arraybackslash}p{0.21\linewidth}>{\raggedright\arraybackslash}p{0.24\linewidth}>{\raggedright\arraybackslash}p{0.47\linewidth}@{}}
    \toprule
    Workload & Ranking & Reason \\
    \midrule
    data in cache & SPED $>$ Flash & Flash tests for memory presence
      unnecessarily \\
    & SPED, Flash $>$ MT, MP & synchronization and context-switching
      overhead \\
    \midrule
    disk-heavy & Flash $>$ SPED & SPED blocks, lacking asynchronous I/O \\
    & Flash $>$ MT, MP & Flash is more memory-efficient and switches
      contexts less \\
    \bottomrule
  \end{tabular}
\end{table}

\begin{description}
  \item[Single file (best case)] With the synthetic workload every request
    is for the same file, which is always cached.  The bandwidth, $n$
    requests times the file size divided by the time, was measured for file
    sizes up to \SI{200}{\kilo\byte}.  All servers show a similar trend.
    SPED performs best.  Flash comes close, but it pays for the
    \texttt{mincore} test on every request, which here is never needed.
    Zeus shows an anomaly at some file sizes, attributed to data not aligned
    for DMA.  MT and MP are slower because of synchronization and context
    switching, and Apache, which lacks the optimizations, is slowest.
  \item[Owlnet trace] The data set is small and mostly fits in the cache,
    so the trends resemble the best case.  Occasionally, however, a request
    needs the disk.  SPED then blocks entirely, while the helpers of Flash
    keep it serving other requests.
  \item[CS trace] The data set is larger, and many requests need disk I/O.
    SPED performs worst, because it lacks asynchronous I/O.  MT is better
    than MP, thanks to its smaller memory footprint and cheaper
    synchronization.  Flash performs best: its small footprint leaves more
    memory for caching, so fewer requests need blocking I/O, and it needs no
    synchronization.
  \item[Optimizations] Measurements of Flash with various combinations of
    the caches of path names, response headers and mapped files show that
    each optimization contributes substantially.  Apache would have
    benefited from them too.
\end{description}

\begin{keypoint}
  When the data are in memory, the event-driven servers win, and the extra
  test of Flash costs a little.  When the workload needs the disk, the
  helpers of Flash avoid blocking, and its small memory footprint leaves
  more memory for caching, which reduces disk I/O.  Flash performs well
  across workloads rather than best for each.
\end{keypoint}

\needspace{6\baselineskip}
\section{Designing and running experiments}
\label{sec:l06-design}

An experiment is worth running only if it is relevant: it must support a
statement about a solution that others believe in and care
about.\source{P2L5, advice on designing experiments, advice on running
experiments.}

\subsection{Relevance and metrics}

Relevance depends on the stakeholders.  For a web server, clients care about
response time, while the operator of the server cares about throughput and
costs.  A new design that improves both response time and throughput is
clearly valuable; one that improves one metric without degrading the other
is still an improvement; one that improves one metric at the expense of the
other may be useful to some stakeholders, provided the degraded metric
remains acceptable.  The goals thus determine which metrics to measure and
how to configure the experiments.

Metrics should be chosen accordingly.  Standard metrics reach a broader
audience, because readers know how to interpret and compare them; for the
same reason, experiments often use a \term{benchmark}[ベンチマーク], a
standardized workload with rules for running it and reporting the results
that a community has agreed upon.  Beyond that, the metrics should answer
the questions of who cares about the result, what they care about and why:
client performance calls for response time or the number of requests that
time out, and operator costs for throughput and costs.

\subsection{Configuration space}

The behaviour of a system depends on many factors, which together form its
configuration space:
\begin{itemize}
  \item \textbf{system resources}: hardware, such as the CPUs and memory,
    and software, such as the number of threads or the sizes of queues;
  \item \textbf{workload}: for a web server, the request rate, the number of
    concurrent requests, the file sizes and the access pattern.
\end{itemize}
Since the whole space cannot be explored, an experiment chooses a subset of
the configuration parameters, picks a range of values for each parameter
that is varied, and uses a relevant workload.  The best-case and worst-case
scenarios should be included.  The combinations of factors should be
useful: many combinations merely repeat the same point.

Two rules govern comparisons.  First, compare \emph{apples to apples}:
change one factor at a time.  If one server is measured with a light
workload on a large machine and another with a heavy workload on a small
machine, no conclusion about the servers is possible, because the difference
may come from the workload, from the resources or from the servers.
Second, compare against a \term{baseline}[ベースライン]: a reference system
whose performance is understood, such as the state of the art, the most
common practice, or an ideal best-case or worst-case scenario.  A claim
that a design is better is a claim that it is better than a baseline.

\subsection{Running experiments}

Once the experiments are designed, each test case is run many times, the
metrics are computed over the runs, for example as averages together with
their spread, and the results are represented in a form, typically a graph,
that makes the point clearly.  Finally, conclusions must be drawn:
measurements that are not interpreted with respect to the goals of the
experiment do not support any statement.

\begin{example}
  A team wants to show that a new AMPED server reduces the 95th-percentile
  response time of a photo-sharing site without lowering throughput.  The
  hardware is fixed at one server machine, and the workload is a trace of
  the site replayed at increasing request rates of 1\,000, 2\,000, 4\,000,
  8\,000 and 16\,000 requests per second.  A best case, with every file in
  memory, and a worst case, with a data set four times larger than memory,
  bound the range.  The baseline is the multithreaded server currently in
  use.  With two servers, five rates, two cases and five repetitions each,
  the study consists of $2 \cdot 5 \cdot 2 \cdot 5 = 100$ runs, and it
  reports for every configuration the mean of the percentile and of the
  throughput together with their spread.
\end{example}

\needspace{11\baselineskip}
\begin{keypoint}[Lesson summary]
  Which design is better depends on the metric and on the workload.  A
  concurrent server can use processes (MP), threads (MT) or one
  event-driven thread (SPED); helpers turn blocking operations into events
  when asynchronous I/O is unavailable (AMPED, AMTED).  Flash, an AMPED
  server, comes close to SPED when data are cached and outperforms all
  alternatives when the workload needs the disk; Apache combines processes
  with dynamic thread pools.  Sound experiments choose relevant metrics and
  workloads, vary one factor at a time, compare against a baseline, repeat
  runs and draw conclusions.
\end{keypoint}

\end{document}
